\documentclass[UTF8,fontset=none,zihao=-4,a4paper]{ctexart}
\usepackage{geometry}
\geometry{margin=2.4cm}
\usepackage{xeCJK}
\usepackage{fontspec}
\usepackage{tabularx}
\usepackage{booktabs}
\usepackage{longtable}
\usepackage{enumitem}
\usepackage{xcolor}
\usepackage[most]{tcolorbox}
\usepackage{titlesec}
\usepackage{fancyhdr}
\usepackage{hyperref}
\hypersetup{colorlinks=true,linkcolor=black,urlcolor=black,
  pdftitle={黄裳行旅散文集 重编方案},pdfauthor={重编项目}}

% --- CJK fonts (Linux: AR PL + WenQuanYi) ---
\setCJKmainfont{AR PL SungtiL GB}
\setCJKsansfont{WenQuanYi Zen Hei}
\setCJKmonofont{WenQuanYi Zen Hei Mono}
\setCJKfamilyfont{zhkai}{AR PL KaitiM GB}
\newcommand{\kai}{\CJKfamily{zhkai}}
% fallback family for rare CJK glyphs absent from the Sung body font
\setCJKfamilyfont{zhfb}{WenQuanYi Zen Hei}
\newcommand{\fb}{\CJKfamily{zhfb}}
\setmonofont{DejaVu Sans Mono}[Scale=0.9]

% --- route symbols missing from the Latin font through WenQuanYi ---
\usepackage{newunicodechar}
\newfontfamily\symbolfont{WenQuanYi Zen Hei}
\newunicodechar{★}{{\symbolfont\char"2605}}
\newunicodechar{○}{{\symbolfont\char"25CB}}
\newunicodechar{●}{{\symbolfont\char"25CF}}
\newunicodechar{·}{{\symbolfont\char"00B7}}
\newunicodechar{↔}{{\symbolfont\char"2194}}
\newunicodechar{→}{{\symbolfont\char"2192}}
\newunicodechar{←}{{\symbolfont\char"2190}}
\newunicodechar{≤}{{\symbolfont\char"2264}}
\newunicodechar{≥}{{\symbolfont\char"2265}}
\newunicodechar{–}{{\symbolfont\char"2013}}
\newunicodechar{※}{{\symbolfont\char"203B}}

\linespread{1.25}
\setlength{\parindent}{2em}
\setlist{itemsep=2pt,topsep=3pt,parsep=0pt}

% headings in sans (黑体)
\titleformat*{\section}{\Large\sffamily\bfseries}
\titleformat*{\subsection}{\large\sffamily\bfseries}
\titleformat*{\subsubsection}{\normalsize\sffamily\bfseries}

\pagestyle{fancy}
\fancyhf{}
\fancyhead[L]{\small\sffamily 黄裳行旅散文集 · 重编方案}
\fancyhead[R]{\small\thepage}
\renewcommand{\headrulewidth}{0.3pt}

\begin{document}
\sloppy



\begin{titlepage}
\centering
\vspace*{3cm}
{\sffamily\bfseries\fontsize{30}{36}\selectfont 黄裳行旅散文集}\\[6pt]
{\kai\Large 重编方案 · 篇目索引与编辑框架}\\[2.2cm]
{\large 一部新选本的结构、选目、导读与注释设计}\\[0.4cm]
{\large（编选研究工作稿）}\\[3cm]
{\large 黄裳（容鼎昌，1919--2012）行旅散文}\\[0.3cm]
{\large 以文史脉络重构 · 八辑结构}\\[2.5cm]
{\normalsize 2026 年 6 月}\\
\vfill
\begin{tcolorbox}[colback=gray!6,colframe=gray!55,boxrule=0.6pt,
  arc=2pt,left=10pt,right=10pt,top=8pt,bottom=8pt,width=0.92\linewidth]
\small
\textbf{版权与编例说明}\quad 黄裳作品仍在版权保护期内（作者 2012 年逝世，
存续至 2062 年）。本文档为\textbf{编选研究工作稿}，依法\textbf{不收录受版权保护的作品全文}；
所录为篇目索引、\textbf{原创内容摘要（≤200 字，非原文复制）}、版本来源、主题与地点标签、
行旅年表、辑目结构、导读与注释框架。必要引用只取短句并注出处；网络流传节选不作可靠底本。
正式出版前须取得版权方授权，并以原书或《黄裳集》核校全部篇目、年份与文字。
\end{tcolorbox}
\end{titlepage}


\tableofcontents

\clearpage

\part{编选方案}

\section{总序草稿（Volume Preface — Draft）}


\begin{quote}\small\itshape
草稿，待全书定目后修订。文中史实凡标〔待核〕者须核定后落定；不得在正式稿中保留未核之确指。
\end{quote}

\medskip\hrule\medskip

\subsection{拟题：行旅即读史——重编黄裳行旅散文集序}


黄裳（1919—2012）一生走过的路，多半是被史书与旧籍先走过一遍的路。他到一处，先看见的往往不是风景，而是风景背后的人与事：秦淮的桨声里有桃花扇的旧影，鸡鸣寺的台阶上叠着几代登临者的脚迹，琉璃厂的书肆是另一种意义上的古战场。于是他的游记从来不只是游记——\textbf{它是由地点触发的文史随笔}：风景、旧迹、掌故、书籍、戏曲、人物与历史兴亡之感，在他笔下交织成一种独有的纹理。


这部选本据此重编。我们不沿用旧书的书名与出版年代为序，而把散落在《锦帆集》《锦帆集外》《关于美国兵》《白门秋柳》《金陵五记》《花步集》《晚春的行旅》以及各种书话集中的行旅文字，按\textbf{时间、地点、场景、主题、文史脉络}重新分门别类，编为八辑。


第一辑「蜀道与战时中国」，是一个青年在抗战后方的行脚——成都、重庆、昆明，山一程水一程，他由旅行者渐渐成为历史的见证者。第二辑「美国兵与异国战争经验」，记下战时中国与一支外来军队相遇的现场〔战时任美军译员，确定〕，重点不在猎奇，而在记录。第三辑「金陵旧梦」是全书的心脏：南京的兴废、秦淮的烟月、老虎桥畔的故人〔指探看周作人事，待核〕，黄裳写得最厚也最深。第四辑「江南行旅」展开苏扬杭沪的文化地理；第五辑「北平与旧都余韵」把行旅与「书旅」打通；第六辑「山川、古迹与历史残影」让名胜废墟成为追问历史的入口；第七辑「书肆、旧书与纸上行旅」呈现他「读书即旅行，访书即考古」的一面；第八辑「人物途中」，以旅途所遇的作家、艺人、藏书家收束全书。


我们尤其留意一种结构：\textbf{同一个地方，他在不同年代写过不止一次}。鸡鸣寺如此，秦淮如此，南京整座城更如此。把这些「同地重访」之作并置编排，便能看见三四十年间一个人的眼光如何随时代变迁——这是单行本各自为政时看不到的纵深，也是本选本想要呈现的东西。


最后须说明编例。黄裳作品仍在版权保护期内，本书在筹备阶段建立的是篇目索引、版本来源、内容摘要与注释框架；凡引用只取短句并注出处，绝不大段移录。所有篇名、年份、出处皆经核校，凡一时不能确定者宁可存疑标「待考」，不敢臆造一字。\textbf{愿这部选本既能供普通读者卧游，也经得起研究者复核。}


〔署名／日期待定〕


\medskip\hrule\medskip

\subsubsection{修订清单（编辑自用）}

\begin{itemize}
\item 待全书定目后，补「全书共八辑、收文 N 篇」的总数。
\item 「老虎桥」「美军译员」等具体史实落定前保留〔待核〕。
\item 视出版形态决定是否加「凡例」单列（注释体例、版本依据、繁简、标点）。
\end{itemize}


\section{选编原则（Editorial Principles）}


本选本不是旧书重印，而是一次\textbf{以文史脉络重构}的专题编选。以下原则统摄选目、分辑、注释与编排。


\medskip\hrule\medskip

\subsection{一、根本判断：黄裳游记 = 「地点触发的文史随笔」}


黄裳写一地，往往由眼前风景/旧迹切入，旋即转入掌故、史事、书籍、戏曲与人物，落点常在\textbf{历史兴亡之感}。因此：


\begin{itemize}
\item \textbf{不以「到此一游」式风景文为标准}，而以「能否打开一段文史」为取舍。
\item 同一地点的不同时期书写（如鸡鸣寺1946/1979）视为\textbf{结构资源}，刻意并置以显时间纵深。
\item 「读书即旅行，访书即考古」是其独特面向，书话中带行踪者\textbf{等同游记}对待（第七辑）。
\end{itemize}


\subsection{二、选目标准（四档，对应数据库 \texttt{inclusion} 字段）}


\begin{enumerate}
\item \textbf{必收}：文史密度高、代表黄裳方法、具史料或文学双重价值、或为名篇定调之作。
\begin{itemize}
\item 例：白门秋柳、金陵杂记（整组/精选）、老虎桥边看「知堂」、秦淮拾梦记、解放后看江南、关于「翻译官」、西泠访书记、琉璃厂。
\end{itemize}
\item \textbf{可选}：题材重要但时评性强、或与必收篇重复地缘/主题，视全书容量取舍。
\begin{itemize}
\item 例：司徒雷登夹着皮包走了以后、采石·当涂·青山、绝代的散文家——张宗子。
\end{itemize}
\item \textbf{资料库（备而不选）}：属同一原书、保留在篇目库以求完整，但偏离行旅文史主线（如战地后勤杂记、序跋），正文一般不收，必要时仅摘短句入导读。
\begin{itemize}
\item 例：《关于美国兵》中的「咖啡与战斗力」「汽车团」「军中文化」「叙言」等。
\end{itemize}
\item \textbf{不宜收}：纯时政评论、纯版本学书话（无行踪）、纯剧评，与「行旅文史」主旨偏离者。判例见数据库 \texttt{not\_recommended}，理由须注明。
\begin{itemize}
\item 例：谈「善本」「题跋」「集部」「禁书」、古书的作伪（纯书话）；《旧戏新谈》中的戏目评谈（纯剧评）。
\end{itemize}
\end{enumerate}


\begin{quote}\small\itshape
\texttt{inclusion} 取值：必收 / 可选 / 资料库；「不宜收」者不进 essays 表，单列于 \texttt{not\_recommended}。三者合计构成"正文规模"的调节阀——核心辑（第三辑）可多必收，整本书的必收总量须与目标篇幅匹配。
\end{quote}

\subsection{三、分辑逻辑（八辑）}


按\textbf{主题—地缘—时间}三维交织分辑（详见 \texttt{outline/new\_anthology\_outline.md}）：

\begin{enumerate}
\item 蜀道与战时中国（西南行旅 / 抗战后方）
\item 美国兵与异国战争经验（战时中外军队接触现场）
\item 金陵旧梦（南京，\textbf{全书核心}）
\item 江南行旅（苏扬杭沪周边 / 水路 / 园林）
\item 北平与旧都余韵（旧城 / 琉璃厂 / 剧场 / 访书）
\item 山川、古迹与历史残影（不限地域，由名胜废墟入史）
\item 书肆、旧书与纸上行旅（带行踪的书话）
\item 人物途中（旅途所遇的作家、艺人、藏书家、戏曲人物）
\end{enumerate}


\textbf{跨辑篇目处理}：一篇可属多辑者（如金陵的访书文），以\textbf{地缘完整性优先}入主辑，于他辑导读中「互见」提示，数据库 \texttt{editorial\_note} 记录双栖关系，避免重复收录。


\subsection{四、编排次序}


\begin{itemize}
\item 辑内\textbf{先编年、后地点}：同辑按写作年代排，凸显作者眼光的成长与历史变迁。
\item 「同地重访」诸篇\textbf{贴近编排}，形成对读。
\item 每辑前置\textbf{导读}（见 \texttt{section\_introductions.md}），点明该辑文史脉络与代表篇。
\end{itemize}


\subsection{五、注释框架}


每篇拟配三类注（繁简按出版形态定）：

\begin{enumerate}
\item \textbf{地理/史迹注}：地名今昔、建置沿革、所涉古迹。
\item \textbf{掌故/人物注}：所引史事、人物（如柳如是、马湘兰、周作人、司徒雷登）简注。
\item \textbf{版本/异文注}：初刊与结集本差异、可考的写作背景与时间。
\end{enumerate}
\begin{itemize}
\item 注释只作背景说明，\textbf{不复制原文}；引用原句限短句并标页。
\end{itemize}


\subsection{六、版权与诚信底线}


\begin{itemize}
\item 全程\textbf{不大段复制受版权保护全文}；摘要为原创转述（≤200字）。
\item 每条事实带 \texttt{confidence} 与 \texttt{source\_reference}；\textbf{查不到即标待核，绝不臆造篇名/年份/出处}。
\item 网络流传节选不作可靠底本，仅作线索；最终以\textbf{原书或《黄裳集》}核校。
\item 正式出版前须取得版权方授权。
\end{itemize}


\subsection{七、独特性诉求（本选本要做到的）}


\begin{enumerate}
\item \textbf{结构创新}：以文史脉络而非书名/年代分辑，首次系统呈现「地点—史事—人物—书—戏」的交织。
\item \textbf{时间纵深}：用「同地重访」对读，呈现1940s—1980s眼光变迁。
\item \textbf{打通书旅与行旅}：把书话中的访书行踪正式纳入「行旅」范畴（第七辑）。
\item \textbf{学术配套}：编年篇目表、地点/主题多维标签、导读与三类注释，超越普通散文选。
\item \textbf{诚信标注}：全程「确定/待核/推测」透明化，可供研究者复核与续补。
\end{enumerate}


\section{《黄裳行旅散文集》新目录方案（初稿）}


\begin{quote}\small\itshape
八辑结构。\textbf{以文史脉络重构，非按原书/年代堆砌。} 篇目均标「确定/待核/推测」；标★为「必收」候选，标○为「可选」候选。多数篇目尚待逐书录入，本稿为可扩展骨架。
\end{quote}

\medskip\hrule\medskip

\subsection{总体设计}


\begin{itemize}
\item \textbf{核心辑}：第三辑「金陵旧梦」（黄裳南京书写最厚，史料与文学俱重）。
\item \textbf{开合}：以第一辑（青年战时行旅，从旅行者到见证者）起，以第八辑（人物途中）收。
\item \textbf{辑内次序}：先编年、后地点；「同地重访」诸篇贴近编排。
\end{itemize}


\medskip\hrule\medskip

\subsection{第一辑　蜀道与战时中国}

\textbf{范围}：《锦帆集》《锦帆集外》中与蜀道、成都、重庆、昆明、桂林、贵阳、西南行旅、抗战后方有关者。 \textbf{主旨}：呈现年轻黄裳如何从旅行者成为历史见证者。


\begin{itemize}
\item ★ 宝鸡——广元（《锦帆集》，蜀道，呼应辑名）
\item 成都散记（《锦帆集》，蜀中）
\item 过徐州（《锦帆集》，旅途串场）
\item 江上杂记（《锦帆集》，江河；江段待核，或入第四辑）
\item 昆明杂记（待核，西南）
\item 茶馆（待核，后方世情）
\item 森林·雨季·山头人（待核，地点待定，或入第二辑）
\item 〔待补：重庆/桂林/贵阳诸篇——据《锦帆集》《锦帆集外》原书续录〕
\end{itemize}


\subsection{第二辑　美国兵与异国战争经验}

\textbf{范围}：《关于美国兵》（上海出版公司，1947，周报丛书第一辑）全书。 \textbf{主旨}：战时中国与外来军队接触的\textbf{现场记录}，非异域猎奇。 \textbf{已得完整目录（15篇）；作为"行旅散文集"的一辑，不整辑照收，按下表分级精选（资料库篇仅存目不入正文）：}


\begin{center}
\footnotesize
\begin{tabularx}{\linewidth}{@{}XX@{}}
\toprule
\textbf{篇名} & \textbf{取舍} \\
\midrule
他们怎样看中国 & ★必收 \\
中国将军怎样应付美国兵 & ★必收 \\
美国兵与女人 & ★必收 \\
为美国兵活着的人们（社会批判） & ★必收 \\
关于"翻译官"（作者自况，宜作辑末收束） & ★必收 \\
人种·职业·人性的大集合 & ○可选 \\
几个人物 & ○可选 \\
前线景象（疑涉滇西/缅北，待核） & ○可选 \\
种种惊异 & ○可选 \\
叙言（序） & 资料库（可摘句入导读） \\
"中国通" & 资料库 \\
"伟大"的S.O.S. & 资料库 \\
咖啡与战斗力 & 资料库 \\
军中文化 & 资料库 \\
汽车团（疑滇缅公路，待核） & 资料库 \\
\bottomrule
\end{tabularx}
\end{center}


\begin{quote}\small\itshape
拟正文规模：必收5 ＋ 可选4，约 5—9 篇；其余6篇偏战地杂记/序，留资料库备查。叙言可摘短句入第二辑导读，保留"以译者眼看美军"的首尾呼应。
\end{quote}

\subsection{第三辑　金陵旧梦　【全书核心】}

\textbf{范围}：《白门秋柳》《金陵五记》中南京、秦淮、白下、鸡鸣寺、下关、旧都记忆诸作。 \textbf{主旨}：以南京为中心，演示「一地—一史—一人」的文史交织，并以同地重访见时间纵深。


按编年（据《金陵五记》目录，均待核）：

\begin{itemize}
\item ★ 白门秋柳（1942）— 辑首定调
\item 旅京随笔（1946）
\item 鸡鸣寺（1946）↔ 重过鸡鸣寺（1979）〔同地重访并置〕
\item 关于「泽存书库」（1946，亦关书话）
\item 访「盔山精舍」（1946）
\item ★ 老虎桥边看「知堂」（1946，记探周作人）
\item ★ 金陵杂记（1947，整组或精选；含柳如是、马湘兰、桃花扇底看南朝、燕子矶、莫愁湖等约二十则）
\item ★ 解放后看江南（1949）
\item ○ 司徒雷登夹着皮包走了以后（1949）
\item ○ 不再做花瓶（1949）
\item 白下书简（1979）
\item ★ 秦淮拾梦记（1979）
\item 南唐二陵（1979，亦可入第六辑）
\item 〔待补：《白门秋柳》专书篇目〕
\end{itemize}


\subsection{第四辑　江南行旅}

\textbf{范围}：苏州、扬州、杭州、上海周边、江南水路、园林、书肆、旧迹。 \textbf{主旨}：突出江南文化地理。


\begin{itemize}
\item 采石·当涂·青山（待核，长江沿线／李白史迹）
\item 杭州杂记（《花步集》，杭州；篇名/归属推测，待核）
\item 海滨消夏记（《银鱼集》，海滨；地点待核，或入第六辑）
\item 江上杂记（与第一辑互见，视江段定归属）
\item 〔待补：苏/扬/杭/沪周边诸篇——多在《花步集》《晚春的行旅》及《黄裳集》各卷。\textbf{注：这些书的逐篇目录在豆瓣/微信读书/中图网等站，均被出口代理按策略拦截（403），WebSearch 也无法还原；须以原书或国图/上图/南图 OPAC 补录。}〕
\end{itemize}


\subsection{第五辑　北平与旧都余韵}

\textbf{范围}：北京、琉璃厂、旧城、剧场、访书、掌故。 \textbf{主旨}：把「行旅」与「书旅」打通。


\begin{itemize}
\item ★ 琉璃厂（见《黄裳自选集》第四辑；北京书肆古玩街，本辑代表篇）
\item 〔待补：旧城/剧场/访书诸篇——从《榆下说书》《翠墨集》及《旧戏新谈》采辑。本辑目前最弱，须优先补采〕
\end{itemize}


\subsection{第六辑　山川、古迹与历史残影}

\textbf{范围}：不限地域，凡由名胜、古迹、寺观、城墙、废墟进入历史追问者。 \textbf{主旨}：历史残影中的兴亡之思。


\begin{itemize}
\item 南唐二陵（1979，与第三辑互见）
\item 〔待补：各集中怀古/访古诸篇〕
\end{itemize}


\subsection{第七辑　书肆、旧书与纸上行旅}

\textbf{范围}：书话中带地点、行踪、访书经历的文章。 \textbf{主旨}：「读书即旅行，访书即考古」。


\begin{itemize}
\item ★ 西泠访书记（《榆下说书》，杭州西泠访书；本辑代表篇，可与第四辑江南互见）
\item 关于「泽存书库」（与第三辑互见，南京藏书旧事）
\item 琉璃厂（与第五辑互见）
\item 〔待补：《榆下说书》《银鱼集》《翠墨集》《珠还记幸》中带行踪的访书文。注意：谈"善本"/题跋/集部/禁书等纯书话无行踪者\textbf{不入本选本}（见数据库 not\_recommended）〕
\end{itemize}


\subsection{第八辑　人物途中}

\textbf{范围}：旅途遇见的作家、艺人、朋友、藏书家、戏曲人物等记人文字。 \textbf{主旨}：以人物收束全书。


\begin{itemize}
\item 老虎桥边看「知堂」（与第三辑互见，记周作人）
\item 余淡心与金陵（《银鱼集》，记余怀＋秦淮；与第三辑互见，「以人写地」典范）
\item ○ 绝代的散文家——张宗子（《银鱼集》，记张岱；偏文学评论，列可选）
\item 〔待补：记艺人/藏书家/戏曲人物诸篇——见《银鱼集》《珠还记幸》《来燕榭文存》《旧戏新谈》〕
\end{itemize}


\medskip\hrule\medskip

\subsection{附录（拟）}

\begin{itemize}
\item 《锦帆集》后记（自述书名取义与成书缘起，待核）
\item 黄裳行旅年表（据已核篇目编年生成）
\item 篇目索引（地点 / 主题 / 时间多维）
\item 版本与异文校记
\end{itemize}


\medskip\hrule\medskip

\subsection{体量与平衡（编辑提示）}

\begin{itemize}
\item 第三辑最厚（核心），第一、四辑次之；第五辑目前最弱，须优先补采，否则考虑与第六辑合并。
\item 跨辑篇以地缘完整性优先入主辑，他辑「互见」，\textbf{严防重复收录}（去重规则见 \texttt{workflow/research\_workflow.md}）。
\end{itemize}


\section{各辑导读草稿（Section Introductions — Draft）}


\begin{quote}\small\itshape
每辑导读 200—400 字骨架，待篇目落定后补具体篇名与篇数。凡具体史实标〔待核〕者须核定。导读只作引导与背景，不复制原文。
\end{quote}

\medskip\hrule\medskip

\subsection{第一辑　蜀道与战时中国}


抗战把一个读书的青年推上了路。成都、重庆、昆明，是后方，也是另一种意义上的前线。这一辑收黄裳早年（约1944—1946）西南行旅与战时见闻之作〔多出《锦帆集》《锦帆集外》，原书篇目待核〕。可注意的是，他写茶馆、写山路、写雨季里的边地人群，眼光已不止于游观——后方社会的世情、战时中国的体温，都在这些速写里。读这一辑，是看一个旅行者怎样一步步变成历史的见证者。


\textbf{代表篇（候选）}：宝鸡——广元（蜀道）、成都散记、过徐州、江上杂记（《锦帆集》）；昆明杂记、茶馆〔均待核〕。


\medskip\hrule\medskip

\subsection{第二辑　美国兵与异国战争经验}


1944年起黄裳被征调任美军译员〔确定〕，得以近距离观察一支外来军队如何嵌入战时中国。这一辑以《关于美国兵》为主体〔全书篇目待录入〕。它的价值不在异域猎奇，而在现场：语言不通处的尴尬、军队与地方社会的摩擦、战争经济下的人际关系。这是中国现代散文里不多见的、由当事人写下的「中外军队接触现场」。


\textbf{代表篇（候选）}：关于"翻译官"（作者自况）、美国兵与女人、为美国兵活着的人们、他们怎样看中国〔《关于美国兵》全15篇已得目录，均待核〕。


\medskip\hrule\medskip

\subsection{第三辑　金陵旧梦　【全书核心】}


南京是黄裳写得最久、最深的城。从1942年的《白门秋柳》，到1946—1947年战后重访的密集书写，到1949年时代转折的现场，再到1979年的重游追忆——同一座城，他写了将近四十年〔编年据《金陵五记》目录，待核〕。这一辑依此编年，让秦淮、鸡鸣寺、莫愁湖、燕子矶在不同年代反复出现，也让柳如是、马湘兰、桃花扇底的南朝与老虎桥畔的故人一一登场。把「鸡鸣寺（1946）」与「重过鸡鸣寺（1979）」这样的同地重访并置，读者便能看见一个人的眼光与一座城的命运如何彼此映照。这是全书的心脏。


\textbf{代表篇（候选）}：白门秋柳、金陵杂记、老虎桥边看「知堂」、秦淮拾梦记、解放后看江南〔均待核〕。


\medskip\hrule\medskip

\subsection{第四辑　江南行旅}


出了南京，便是更广的江南——苏州、扬州、杭州，以及上海周边的水路与园林。这一辑突出江南的文化地理：一座园林背后的造园者，一条水路连起的市镇与书肆，一处旧迹牵出的诗人墓与传说。黄裳写江南，仍是「即地读史」的路数〔篇目多在《花步集》《晚春的行旅》及《黄裳集》各卷，待录入〕。


\textbf{代表篇（候选）}：采石·当涂·青山〔待核，长江沿线、李白史迹〕。


\medskip\hrule\medskip

\subsection{第五辑　北平与旧都余韵}


这一辑把「行旅」与「书旅」打通。北京的旧城、剧场与琉璃厂，在黄裳笔下既是地理，也是书史与戏史的现场——访书是另一种登临，旧书肆是另一种古迹。〔本辑目前仅《琉璃厂》一篇已得线索，须从书话/记戏诸集（《榆下说书》《翠墨集》《旧戏新谈》）补采；若篇目不足，考虑与第六辑统筹。〕


\textbf{代表篇（候选）}：琉璃厂〔见《黄裳自选集》，待核〕；余待补。


\medskip\hrule\medskip

\subsection{第六辑　山川、古迹与历史残影}


不限地域。凡是从一座城墙、一处寺观、一道废墟出发，最终走向历史追问的文字，归入此辑。黄裳面对古迹时常有的，不是凭吊的伤感，而是考据式的好奇与兴亡之际的清醒。


\textbf{代表篇（候选）}：南唐二陵〔与第三辑互见〕；其余待补。


\medskip\hrule\medskip

\subsection{第七辑　书肆、旧书与纸上行旅}


黄裳是著名藏书家，他的书话里藏着大量行踪：哪一年在哪座城的哪家书肆，遇见哪一部书，背后又牵出怎样的版本与史事。这一辑专收这类「带脚印的书话」，呈现他「读书即旅行，访书即考古」的一面〔来源：《榆下说书》《银鱼集》《翠墨集》《珠还记幸》等，部分已得目录〕。须注意区分：带行踪的访书文（收）与纯版本学书话（如谈"善本"、谈"题跋"，不收）。


\textbf{代表篇（候选）}：西泠访书记〔杭州西泠访书，《榆下说书》，待核〕、琉璃厂〔与第五辑互见〕、关于「泽存书库」〔与第三辑互见〕。


\medskip\hrule\medskip

\subsection{第八辑　人物途中}


路上遇见的人，往往比风景更难忘。作家、艺人、藏书家、戏曲中的旧角色，都在黄裳的旅途里留下过身影。这一辑以记人之文收束全书——让人物，而不是风景，成为最后的落点。


\textbf{代表篇（候选）}：老虎桥边看「知堂」〔记周作人，与第三辑互见〕、余淡心与金陵〔记余怀＋秦淮，《银鱼集》，"以人写地"典范〕、绝代的散文家——张宗子〔记张岱，偏评论，可选〕。


\medskip\hrule\medskip

\subsubsection{编辑提示}

\begin{itemize}
\item 各辑导读须在篇目定稿后补：本辑收文 N 篇、起讫年代、与他辑互见篇说明。
\item 凡〔待核〕史实落定前不得在正式稿中作确指。
\end{itemize}


\part{篇目、年表与资料}

\section{篇目总表（含原创内容摘要）}
\noindent\small 全表 59 篇，按拟定辑次排列。每篇标注〔编号｜取舍｜置信〕。
摘要均为原创转述（≤200 字），\textbf{非原文复制}；篇名、年代、出处凡未坐实者标「待核 / 推测」。\normalsize
\medskip
\subsection{第1辑　蜀道与战时中国（9 篇）}
\noindent\textbf{《宝鸡——广元》}\,{\small ［E-014｜★必收｜待核］}\\
{\small\itshape 原收入：锦帆集　｜　年代：约1944—1946　｜　地点：宝鸡（陕）—广元（川），蜀道（蜀中）}\\
日记体纪行，记由陕西宝鸡入川至广元一线蜀道行程：穿越秦岭、越剑门、入川盆地，描绘沿途关隘峰壑与战时人流；亦收入《往事如烟》等选集。第一辑蜀道行旅最核心的地理线索篇。
\par\vspace{5pt}
\noindent\textbf{《成都散记》}\,{\small ［E-015｜★必收｜待核］}\\
{\small\itshape 原收入：锦帆集　｜　年代：约1944—1946　｜　地点：成都（蜀中）}\\
1943—1946年任美军译员期间对成都后方生活的散记：写宽街窄巷、市井风物与战时文化氛围；与《茶馆》（E-036）形成互补，共同呈现抗战时期蜀中都市的生活横切面。《锦帆集》蜀道组核心篇之一。
\par\vspace{5pt}
\noindent\textbf{《过徐州》}\,{\small ［E-016｜○可选｜待核］}\\
{\small\itshape 原收入：锦帆集　｜　年代：约1944—1946　｜　地点：徐州（江北）}\\
1942年冬（腊月初旬）离沪赴蜀途中过徐州：见集市陈列布匹铜器，与同行者（含黄宗江、宋淇等）购买徐州风景明信片寄回上海；是《锦帆集》蜀道诸篇序章，战时离乡漂泊之感跃然。
\par\vspace{5pt}
\noindent\textbf{《江上杂记》}\,{\small ［E-017｜○可选｜待核］}\\
{\small\itshape 原收入：锦帆集　｜　年代：约1944—1946　｜　地点：长江（舟行，具体江段待核）（江南）}\\
舟行江上的杂记，记水路见闻与过境感受；取书名意象'锦帆应是到天涯'之境。确切江段（川江或长江下游）待核，归辑随之调整；是《锦帆集》少数以江河场景为主的行旅篇。
\par\vspace{5pt}
\noindent\textbf{《昆明杂记》}\,{\small ［E-018｜★必收｜待核］}\\
{\small\itshape 原收入：锦帆集外（文学丛刊第九集，文化生活出版社，1948；含五部分）　｜　年代：约1944—1946　｜　地点：昆明（西南）}\\
任译员期间在昆明的见闻杂记，分五部分；在昆明购得载永历帝滇缅流亡的晚明野史，三百年后重走同一片土地，遂以鲁迅'抄古书'笔法写历史笔记式杂记，沧桑感浓；文气近汪曾祺。收《锦帆集外》（文学丛刊第九集）。
\par\vspace{5pt}
\noindent\textbf{《森林·雨季·山头人》}\,{\small ［E-019｜○可选｜待核］}\\
{\small\itshape 原收入：锦帆集（待核）　｜　年代：约1944—1946　｜　地点：西南/滇缅印一带（待核）（西南／海外）}\\
战时行役中关于森林、雨季与边地人群的记述，地理指向待核（可能涉滇缅印战区）。具民族志色彩。
\par\vspace{5pt}
\noindent\textbf{《茶馆》}\,{\small ［E-036｜○可选｜待核］}\\
{\small\itshape 原收入：锦帆集（待核）　｜　年代：约1944—1946　｜　地点：西南（成渝一带，待核）（蜀中／西南）}\\
写西南后方各地茶馆的生态差异与阶层混杂：广元/剑阁茶叶粗砺，成都茶馆松弛，重庆茶馆兼作谈判场——警察与挑夫同座、大学生以茶馆为自修室、生意人在此成交。一室之中浓缩战时社会横截面，风俗志价值极高。
\par\vspace{5pt}
\noindent\textbf{《断片》}\,{\small ［E-055｜○可选｜待核］}\\
{\small\itshape 原收入：锦帆集　｜　年代：约1944—1946　｜　地点：（追忆，非行旅地点）（蜀中／西南）}\\
《锦帆集》第二章，题名'断片'；据研究，此篇追忆少年时期的初恋情感，是全书调性最私密的篇目，与其余蜀道行旅诸篇形成对比；作为全集青春记忆的底色，是理解黄裳早年情感世界的罕见文本。
\par\vspace{5pt}
\noindent\textbf{《音尘》}\,{\small ［E-056｜○可选｜推测］}\\
{\small\itshape 原收入：锦帆集　｜　年代：约1944—1946　｜　地点：西南后方（待核）（蜀中／西南）}\\
《锦帆集》第七章，以'音尘'为题（音讯，兼含离别与消逝之叹）；据研究，写战时音讯阻隔与离别之痛，是全书情感基调最深沉的一篇；《音尘集》（1996年辽宁教育版）即以此篇命名，彰显其在黄裳自我认知中的特殊分量。
\par\vspace{5pt}
\subsection{第2辑　美国兵与异国战争经验（16 篇）}
\noindent\textbf{《叙言》}\,{\small ［E-020｜·资料库｜待核］}\\
{\small\itshape 原收入：关于美国兵　｜　年代：约1947　｜　地点：—（—）}\\
《关于美国兵》全书序言，交代写作缘起与立场。导读价值高。
\par\vspace{5pt}
\noindent\textbf{《他们怎样看中国》}\,{\small ［E-021｜★必收｜待核］}\\
{\small\itshape 原收入：关于美国兵　｜　年代：约1944—1947　｜　地点：西南后方（待核）（西南）}\\
记录美军士兵对中国的具体印象：普通士兵称中国人'聪明高尚但缺乏教育'，芝加哥大学毕业军官称中国人'缺乏组织、散漫成群'；以具体人物的实际言论呈现中美文化认知落差，是《关于美国兵》中中美文化观察最集中的篇章。
\par\vspace{5pt}
\noindent\textbf{《“中国通”》}\,{\small ［E-022｜·资料库｜待核］}\\
{\small\itshape 原收入：关于美国兵　｜　年代：约1944—1947　｜　地点：西南后方（待核）（西南）}\\
写美军中自诩'中国通'的外籍军官或专家：以具体人物的语言习惯与认知偏差呈现文化误读与傲慢，讽刺意味浓；与《他们怎样看中国》（E-021）构成中美文化认知专题的两面。
\par\vspace{5pt}
\noindent\textbf{《人种·职业·人性的大集合》}\,{\small ［E-023｜○可选｜待核］}\\
{\small\itshape 原收入：关于美国兵　｜　年代：约1944—1947　｜　地点：西南后方（待核）（西南）}\\
写美军内部各色人种（黑人、白人、拉丁裔等）、职业（军官/士兵/文职）与人性混杂景观；从社会学角度呈现美军作为一个社会横截面的复杂性，群像式观察，史料感强。
\par\vspace{5pt}
\noindent\textbf{《几个人物》}\,{\small ［E-024｜○可选｜待核］}\\
{\small\itshape 原收入：关于美国兵　｜　年代：约1944—1947　｜　地点：西南后方（待核）（西南）}\\
记战时接触的几位美军人物素描：通过具体个人的性格、言行与经历呈现美军群像；人物速写式，文学性强，与第八辑'人物途中'的写法呼应。具体人物待核原文。
\par\vspace{5pt}
\noindent\textbf{《中国将军怎样应付美国兵》}\,{\small ［E-025｜★必收｜待核］}\\
{\small\itshape 原收入：关于美国兵　｜　年代：约1944—1947　｜　地点：西南后方（待核）（西南）}\\
写中国将军应对美军联合指挥的方式：以官场智慧周旋于美方僵化指挥与本方利益之间，现场政治与人情交织；以具体军官行为呈现中美联军权力结构中的张力与博弈。
\par\vspace{5pt}
\noindent\textbf{《美国兵与女人》}\,{\small ［E-026｜★必收｜待核］}\\
{\small\itshape 原收入：关于美国兵　｜　年代：约1944—1947　｜　地点：西南后方（待核）（西南）}\\
写战时美军与中国社会接触中关于女性的一面：美军基地周边的娱乐场所、中美战时婚恋关系与性别风气变迁；纪实性强，是《关于美国兵》最具社会学价值的一章，亦见《黄裳自选集》。
\par\vspace{5pt}
\noindent\textbf{《“伟大”的S.O.S.》}\,{\small ［E-027｜·资料库｜待核］}\\
{\small\itshape 原收入：关于美国兵　｜　年代：约1944—1947　｜　地点：西南后方（待核）（西南）}\\
S.O.S.即美军后勤勤务部（Services of Supply），以反讽笔调写其'伟大'运作：繁复补给程序、物资浪费与官僚机器荒诞；物质批判与战争荒诞感合一，标题引号即反讽所在。
\par\vspace{5pt}
\noindent\textbf{《前线景象》}\,{\small ［E-028｜○可选｜待核］}\\
{\small\itshape 原收入：关于美国兵　｜　年代：约1944—1947　｜　地点：前线（待核，疑滇西/缅北）（西南／海外）}\\
记前线见闻：战壕、废墟、炮火下的日常与死亡气息；具体战区待核（疑滇西或缅北反攻一线），是《关于美国兵》中从后方视角转向真实战场的一章，现场感最强。
\par\vspace{5pt}
\noindent\textbf{《咖啡与战斗力》}\,{\small ［E-029｜·资料库｜待核］}\\
{\small\itshape 原收入：关于美国兵　｜　年代：约1944—1947　｜　地点：西南后方（待核）（西南）}\\
以咖啡这一美军标配物资为切口写后勤与战斗力：充足补给、规律饮食与高昂士气，与中国军队物质匮乏形成对照；借小物件见战争逻辑与国力差异，角度新颖独到。
\par\vspace{5pt}
\noindent\textbf{《种种惊异》}\,{\small ［E-030｜○可选｜待核］}\\
{\small\itshape 原收入：关于美国兵　｜　年代：约1944—1947　｜　地点：西南后方（待核）（西南）}\\
记中美接触中的种种文化惊异：饮食习惯、礼仪观念、价值判断的错位与误解；以具体场景呈现文化落差，宜与《他们怎样看中国》（E-021）对读，共构第二辑文化碰撞主题。
\par\vspace{5pt}
\noindent\textbf{《军中文化》}\,{\small ［E-031｜·资料库｜待核］}\\
{\small\itshape 原收入：关于美国兵　｜　年代：约1944—1947　｜　地点：西南后方（待核）（西南）}\\
写美军中丰富的文化生活与阅读资源：美军随军携带《生活》《时代》《展望》《纽约客》等杂志，驻地自办军中报纸，书架备有奥尼尔戏剧集、狄更斯小说等文学经典；与中国战时文化资源的匮乏形成强烈对照。
\par\vspace{5pt}
\noindent\textbf{《汽车团》}\,{\small ［E-032｜·资料库｜待核］}\\
{\small\itshape 原收入：关于美国兵　｜　年代：约1944—1947　｜　地点：西南后方（待核，疑滇缅公路沿线）（西南）}\\
写美军汽车团沿滇缅公路长途卡车运输的艰险：历经'二十四道拐'险途与西南山区路况，驾驶员以机器维系后方命脉；以具体运输生态呈现抗战后勤战争的真实面貌，是《关于美国兵》中场景感最强的一章。
\par\vspace{5pt}
\noindent\textbf{《为美国兵活着的人们》}\,{\small ［E-033｜★必收｜待核］}\\
{\small\itshape 原收入：关于美国兵　｜　年代：约1944—1947　｜　地点：西南后方（待核）（西南）}\\
写战时围绕美军基地谋生的中国底层人群：街头摊贩、洗衣女工、酒吧女招待、翻译助手……生计与尊严皆系于美军的存在。以底层视角呈现战时经济依附性，社会批判意味浓烈，是《关于美国兵》中阶级观察最深刻的一章。
\par\vspace{5pt}
\noindent\textbf{《关于“翻译官”》}\,{\small ［E-034｜★必收｜待核］}\\
{\small\itshape 原收入：关于美国兵　｜　年代：约1944—1947　｜　地点：西南后方（待核）（西南）}\\
以作者本人译员经历为视角，辨析'翻译官'称呼背后的身份尴尬：外事局称'X级译员'，远征军称'Interpreting Officer'，民间戏称'姨太太''二毛子'，折射中美权力关系在日常语言中的投影；夫子自道，是第二辑最贴近作者战时身份的收束之篇。
\par\vspace{5pt}
\noindent\textbf{《印度小夜曲》}\,{\small ［E-060｜★必收｜待核］}\\
{\small\itshape 原收入：音尘集（辽宁教育出版社，书趣文丛第三辑，1996）　｜　年代：约1944—1946　｜　地点：印度（加尔各答、菩提加雅等）（海外）}\\
黄裳1944—1946年随中美混合部队赴印度期间的行旅长篇，分八章：序曲、雨季、加尔各答、菩提加雅、群莺乱飞、弄蛇者、女人们、森林雨季。以南亚异域风土与战时情境为主要内容，是《关于美国兵》之外战时海外经验的另一维度。
\par\vspace{5pt}
\subsection{第3辑　金陵旧梦（24 篇）}
\noindent\textbf{《白门秋柳》}\,{\small ［E-001｜★必收｜确定］}\\
{\small\itshape 原收入：锦帆集（1946，写于1942；后收入《金陵五记》）　｜　年代：1942　｜　地点：南京（白门）（金陵）}\\
1942年冬黄裳与黄宗江、宋淇（宋希）、周杲良四人离沪赴川，途经汪伪'首都'南京；描绘沦陷南京的衰飒——鸡鸣寺、随园、扫叶楼凋落，书店中《人间味》杂志与'今日的江南的无声的悲哀'；称之为'揉碎了的美'的战时见证。
\par\vspace{5pt}
\noindent\textbf{《旅京随笔》}\,{\small ［E-002｜○可选｜待核］}\\
{\small\itshape 原收入：金陵五记　｜　年代：1946　｜　地点：南京（金陵）}\\
1946年黄裳以记者身份赴南京，目睹国民政府'劫收'乱象冲击民生文化；乘机游览名胜、探访旧日友人与藏书家。'旅京随笔'是《金陵五记》1946年编总称，下含鸡鸣寺、泽存书库、盔山精舍、老虎桥知堂、美人肝共五篇。
\par\vspace{5pt}
\noindent\textbf{《鸡鸣寺》}\,{\small ［E-003｜★必收｜待核］}\\
{\small\itshape 原收入：金陵五记　｜　年代：1946　｜　地点：南京·鸡鸣寺（金陵）}\\
1946年黄裳以记者身份访鸡鸣寺，欲于豁蒙楼写稿，发现楼已被'防空司令部'征用，只得在大殿内喝茶，眺望后湖紫金山；记陈后主'胭脂井'典故（三人同束以绳引出）与南朝兴亡之感。与1979年《重过鸡鸣寺》构成同地重访。
\par\vspace{5pt}
\noindent\textbf{《关于“泽存书库”》}\,{\small ［E-004｜○可选｜待核］}\\
{\small\itshape 原收入：金陵五记　｜　年代：1946　｜　地点：南京（金陵）}\\
记汪伪官员陈群在南京所建'泽存书库'（取《礼记》'手泽存焉'之义）：藏书40万册、善本4400余部，含汪精卫旧藏；战后被接收，善本随蒋携台。访书过程中书史与汪伪政权兴亡交织，史料价值极高。
\par\vspace{5pt}
\noindent\textbf{《访“盔山精舍”》}\,{\small ［E-005｜○可选｜待核］}\\
{\small\itshape 原收入：金陵五记　｜　年代：1946　｜　地点：南京（金陵）}\\
访南京'盔山精舍'的记游兼记人之作。所涉人物与建置待核。
\par\vspace{5pt}
\noindent\textbf{《老虎桥边看“知堂”》}\,{\small ［E-006｜★必收｜待核］}\\
{\small\itshape 原收入：金陵五记　｜　年代：1946　｜　地点：南京·老虎桥监狱（金陵）}\\
1946年黄裳以《文汇报》记者身份赴南京，专赴老虎桥模范监狱探看周作人（知堂）；经延龄巷、国府路找到监狱，申请获准后会面，完成于最终宣判前。长篇报道性质，是战后审奸语境下的人物现场记录，史料价值极高。
\par\vspace{5pt}
\noindent\textbf{《金陵杂记》}\,{\small ［E-007｜★必收｜待核］}\\
{\small\itshape 原收入：金陵五记　｜　年代：1947　｜　地点：南京（多处）（金陵）}\\
1947年系列短札，下含小序及半山寺与谢公墩、石观音寺、周处读书台、绛云书卷美人图、柳如是、快园、随园、梅园、后湖、小虹桥、咏怀堂诗、梅花山、莫愁湖、明太祖与徐达、桃花扇底看南朝、马湘兰、记者生涯、豁蒙楼上看浓春、燕子矶、白鹭洲、鸡鹅巷与裤子裆等二十二则。
\par\vspace{5pt}
\noindent\textbf{《解放后看江南》}\,{\small ［E-008｜★必收｜待核］}\\
{\small\itshape 原收入：金陵五记　｜　年代：1949　｜　地点：江南／南京（江南）}\\
1949年秋南京解放，黄裳亲历人民迎接解放的欢欣；同时目睹满目疮痍、民生凋蔽——劫后时代的两面。历史巨变的现场见证，兼具时代兴奋与直面现实的张力，是《金陵五记》1949年编核心篇。
\par\vspace{5pt}
\noindent\textbf{《司徒雷登夹着皮包走了以后》}\,{\small ［E-009｜○可选｜待核］}\\
{\small\itshape 原收入：金陵五记　｜　年代：1949　｜　地点：南京（金陵）}\\
1949年以美国驻华大使司徒雷登携皮包黯然离华为楔子，写南京政局剧变中美国势力退出与国民政府垮台的亲历观察；与毛泽东同年《别了，司徒雷登》形成历史互文。时评性强，史料价值可观。
\par\vspace{5pt}
\noindent\textbf{《白下书简》}\,{\small ［E-010｜○可选｜待核］}\\
{\small\itshape 原收入：金陵五记　｜　年代：1979　｜　地点：南京（白下）（金陵）}\\
1979年以书简体重写南京，收十一篇；格调低沉，距早年金陵记忆逾三十年，文革后重访兼含追忆旧迹与今昔对照。'白下'为南京历史别称。是《金陵五记》1979年编五类文字之一。
\par\vspace{5pt}
\noindent\textbf{《秦淮拾梦记》}\,{\small ［E-011｜★必收｜待核］}\\
{\small\itshape 原收入：金陵五记　｜　年代：1979　｜　地点：南京·秦淮（金陵）}\\
1979年重游秦淮，以余怀《板桥杂记》为引，追溯秦淮旧影；涉及柳如是、陈圆圆、余淡心等明末清初人物；是1979年南京访古组最具文学性的一篇，网络有节选流传（须以原书核全文）。
\par\vspace{5pt}
\noindent\textbf{《重过鸡鸣寺》}\,{\small ［E-012｜★必收｜待核］}\\
{\small\itshape 原收入：金陵五记　｜　年代：1979　｜　地点：南京·鸡鸣寺（金陵）}\\
约1980年写的重访记，构建三重时空对话：陈后主亡国（胭脂井典故）、十七世纪余怀父子文字（《金陵览古》）、1980年重游；豁蒙楼本为张之洞纪念门生杨锐所建，几度易主；是全书'同地重访'结构最完整的一篇。
\par\vspace{5pt}
\noindent\textbf{《南唐二陵》}\,{\small ［E-013｜○可选｜待核］}\\
{\small\itshape 原收入：金陵五记　｜　年代：1979　｜　地点：南京·南唐二陵（金陵）}\\
1979年访南京祖堂山南唐二陵（先主李{\fb 昪}钦陵、中主李{\fb 璟}顺陵），以词帝李煜亡国为历史背景；追溯南唐短暂王朝（仅39年）的文艺辉煌与政治覆灭，是金陵兴亡感最浓烈的一篇怀古文字。
\par\vspace{5pt}
\noindent\textbf{《「美人肝」》}\,{\small ［E-044｜○可选｜待核］}\\
{\small\itshape 原收入：金陵五记　｜　年代：1946　｜　地点：南京（金陵）}\\
写南京名馔'美人肝'（鸭胰脏）：此菜与松鼠鱼、蛋烧卖、凤尾虾并称'京苏菜四大名菜'，品味于清真老字号马祥兴馆；借饮食掌故钩沉金陵饮食文化传统，文气轻灵，兼具考证趣味与口腹之乐，是《金陵五记》1946年组中风格最轻松的一篇。
\par\vspace{5pt}
\noindent\textbf{《社会科学者眼里的南京》}\,{\small ［E-045｜○可选｜待核］}\\
{\small\itshape 原收入：金陵五记　｜　年代：1949　｜　地点：南京（金陵）}\\
1949年以社会观察者视角审视南京政治、经济与社会变动。内容待核。与同年《解放后看江南》《司徒雷登夹着皮包走了以后》同属1949年南京见证书写。
\par\vspace{5pt}
\noindent\textbf{《不再做花瓶——记中国农业机械公司南京工厂》}\,{\small ［E-046｜○可选｜待核］}\\
{\small\itshape 原收入：金陵五记　｜　年代：1949　｜　地点：南京（金陵）}\\
1949年记南京某工厂（中国农业机械公司）的纪实文字，属解放初期社会面貌观察。内容与文学质量待核。
\par\vspace{5pt}
\noindent\textbf{《白鹭洲公园补记》}\,{\small ［E-047｜○可选｜待核］}\\
{\small\itshape 原收入：金陵五记　｜　年代：1979　｜　地点：南京·白鹭洲（金陵）}\\
1979年对白鹭洲公园的补记。白鹭洲为南京历史园林，有秦淮文化脉络；此篇与《金陵杂记》1947年白鹭洲条构成同地重访对照。内容待核。
\par\vspace{5pt}
\noindent\textbf{《天王府的西花园》}\,{\small ［E-048｜★必收｜待核］}\\
{\small\itshape 原收入：金陵五记　｜　年代：1979　｜　地点：南京·天王府（太平天国旧址）（金陵）}\\
1979年访太平天国天王府遗址之西花园（今南京总统府区域）：以历史层积视角呈现一地三叠——太平天国洪秀全宫室、民国国民政府、当代公园；黄裳追溯晚清历史地图与今昔之变，金陵近代史迹密度最高、时间跨度最宽的一篇。
\par\vspace{5pt}
\noindent\textbf{《梅园（1979版）》}\,{\small ［E-049｜○可选｜待核］}\\
{\small\itshape 原收入：金陵五记　｜　年代：1979　｜　地点：南京·梅园（金陵）}\\
1979年专题写梅园新村（1946年国共谈判中共代表团驻地，周恩来曾居此）：地点今已辟为纪念馆，以亲历者视角追忆三十年前谈判史迹与政治氛围；与《金陵杂记》1947年梅园子目构成时间纵深，重访结构清晰。
\par\vspace{5pt}
\noindent\textbf{《王介甫与金陵》}\,{\small ［E-050｜★必收｜待核］}\\
{\small\itshape 原收入：金陵五记　｜　年代：1979　｜　地点：南京（半山园等王安石相关遗迹）（金陵）}\\
写王安石（字介甫）晚年退居金陵的行迹与精神面貌：'骑驴粗服往来钟山'；追溯其《桂枝香·金陵怀古》被后人称为绝唱（'诸公寄调桂枝香者三十余家，惟王介甫为绝唱'）；政治家与词人两面在金陵山川中交叠，史论与地理感并茂。
\par\vspace{5pt}
\noindent\textbf{《莫愁湖（1979版）》}\,{\small ［E-051｜○可选｜待核］}\\
{\small\itshape 原收入：金陵五记　｜　年代：1979　｜　地点：南京·莫愁湖（金陵）}\\
1979年专题写莫愁湖：钩沉'莫愁女'传说（南朝民歌《莫愁乐》主人公）与'胜棋楼'朱元璋赐名典故；与《金陵杂记》1947年莫愁湖条构成三十年重访对照，展现金陵历史地标在现代的延续与变迁。须核定是否独立成篇。
\par\vspace{5pt}
\noindent\textbf{《石巢园》}\,{\small ［E-052｜○可选｜待核］}\\
{\small\itshape 原收入：金陵五记　｜　年代：1979　｜　地点：南京·石巢园（金陵）}\\
1979年访南京石巢园——明末清初人物阮大铖的宅园（由造园大家计成参与设计）；阮大铖以'库司坊'（民间蔑称'裤子裆'）筑此园，后降清变节；黄裳以园林精美与园主人格的反差呈现历史道德的复杂性，文史批判力度强。
\par\vspace{5pt}
\noindent\textbf{《扫叶楼》}\,{\small ［E-053｜○可选｜待核］}\\
{\small\itshape 原收入：金陵五记　｜　年代：1979　｜　地点：南京·扫叶楼（清凉山）（金陵）}\\
1979年访南京清凉山扫叶楼：楼以清初画家龚贤（1618—1689）旧居兼画室得名，龚贤曾穿僧服作扫叶状以示不事清廷；黄裳追溯金石画史与气节人格，书画史与金陵山川地脉交织，是第三辑中画史掌故最丰富的一篇。
\par\vspace{5pt}
\noindent\textbf{《献花岩》}\,{\small ［E-054｜○可选｜待核］}\\
{\small\itshape 原收入：金陵五记　｜　年代：1979　｜　地点：南京郊外·献花岩（具体地点待核）（金陵）}\\
1979年访献花岩（南京近郊某山川古迹）。具体位置与内容待核，可能涉山川探胜与历史追想。
\par\vspace{5pt}
\subsection{第4辑　江南行旅（3 篇）}
\noindent\textbf{《采石·当涂·青山》}\,{\small ［E-035｜○可选｜待核］}\\
{\small\itshape 原收入：锦帆集（待核）　｜　年代：待核　｜　地点：采石矶、当涂、青山（皖南／长江沿线）（江南）}\\
沿长江访采石矶、当涂、青山李白墓一线：引陆游《入蜀记》勾勒采石矶山水；述李白初赴当涂所赋《望天门山》；追溯范传正将李白迁葬青山始末。山水、史迹、人物三合一，是第四辑江南行旅的典范之作。
\par\vspace{5pt}
\noindent\textbf{《海滨消夏记》}\,{\small ［E-039｜○可选｜待核］}\\
{\small\itshape 原收入：银鱼集（后并入《黄裳集·创作卷V》）　｜　年代：待核　｜　地点：海滨（具体地待核，疑青岛/北戴河一带）（待核）}\\
海滨（疑青岛或北戴河）避暑期间的行旅散记，旁及散文美学漫想（'科学著作也可成美文'），文风轻松闲适；属黄裳较少见的休闲题材，可丰富全书行旅调性的层次，与战时纪行形成对照。
\par\vspace{5pt}
\noindent\textbf{《杭州杂记》}\,{\small ［E-043｜○可选｜待核］}\\
{\small\itshape 原收入：花步集（待核）　｜　年代：待核　｜　地点：杭州（江南）}\\
写杭州（西湖一带）的城市游记，已见于《黄裳散文选集》第一辑记游篇目，篇目真实性较可靠；具体内容与原书归属（疑《花步集》）仍待核，须以原书确认篇名、文字与归集。
\par\vspace{5pt}
\subsection{第5辑　北平与旧都余韵（1 篇）}
\noindent\textbf{《琉璃厂》}\,{\small ［E-038｜★必收｜待核］}\\
{\small\itshape 原收入：待核（见《黄裳自选集》第四辑；疑出书话诸集）　｜　年代：待核　｜　地点：北京·琉璃厂（北平）}\\
写北京琉璃厂旧书古玩街的访书见闻：从清咸丰年间始创的来薰阁，到1925年陈济川接掌后旧书业的领军地位；书肆掌故与北平城市记忆交织，在书话与行旅之间打通'书旅'主题，是第五辑故都题材与第七辑书肆题材的双重核心篇。
\par\vspace{5pt}
\subsection{第7辑　书肆、旧书与纸上行旅（1 篇）}
\noindent\textbf{《西泠访书记》}\,{\small ［E-037｜★必收｜待核］}\\
{\small\itshape 原收入：榆下说书（后并入《黄裳集·创作卷V》）　｜　年代：待核　｜　地点：杭州·西泠（西湖孤山一带）（江南）}\\
记在杭州西湖孤山一带（西泠）寻访旧书铺的经历：书肆散布于园林湖山之间，访书即是行旅，翻检中偶遇珍本；书事、人事与西湖景致交织，是'旅行即读书'主题的最佳演示，亦是第七辑书肆旅行专辑的核心篇。
\par\vspace{5pt}
\subsection{第8辑　人物途中（5 篇）}
\noindent\textbf{《余淡心与金陵》}\,{\small ［E-040｜★必收｜待核］}\\
{\small\itshape 原收入：银鱼集（后并入《黄裳集·创作卷V》）　｜　年代：待核　｜　地点：南京（秦淮）（金陵）}\\
以明末清初文人余怀（字淡心，著《板桥杂记》）游宦金陵的行迹为线索：钩沉其在秦淮河畔的交游与写作背景，追究《板桥杂记》的底本意义；以人物串联金陵文学史，是'以人写地'写法的典范，与E-057互补，建议第三辑/第八辑各取其一。
\par\vspace{5pt}
\noindent\textbf{《绝代的散文家——张宗子》}\,{\small ［E-041｜○可选｜待核］}\\
{\small\itshape 原收入：银鱼集（后并入《黄裳集·创作卷V》）　｜　年代：待核　｜　地点：—（涉绍兴/杭州等张岱行迹）（江南）}\\
以'绝代的散文家'高度评价明末张岱（字宗子）：细析其山水小品的感官精准与故国之思，以《陶庵梦忆》《西湖梦寻》为例；黄裳自身趣味与张岱遥相呼应，是理解其文学观的重要参照。属论人之作，行旅成分弱，列可选。
\par\vspace{5pt}
\noindent\textbf{《关于余淡心》}\,{\small ［E-057｜○可选｜待核］}\\
{\small\itshape 原收入：银鱼集（后并入《黄裳集·创作卷V》）　｜　年代：待核　｜　地点：—（涉南京秦淮等余怀相关地点）（金陵）}\\
关于明末清初文人余怀（字淡心，著《板桥杂记》）的评论性记人文字，与E-040《余淡心与金陵》互见。前者以金陵地点为主，此篇侧重人物评述。内容待核。
\par\vspace{5pt}
\noindent\textbf{《关于吴梅村》}\,{\small ［E-058｜○可选｜待核］}\\
{\small\itshape 原收入：银鱼集（后并入《黄裳集·创作卷V》）　｜　年代：待核　｜　地点：—（涉江南，吴伟业相关行迹）（江南）}\\
写清初诗人吴伟业（字骏公，号梅村）的贰臣处境：以诗文成就与仕清后的自悔心理为核心，追溯其《圆圆曲》背后的政治遭遇与道德煎熬；黄裳对晚明清初气节问题的持续关注在此集中体现，是第八辑人物书写中思想张力最强的一篇。
\par\vspace{5pt}
\noindent\textbf{《咏怀堂》}\,{\small ［E-059｜○可选｜待核］}\\
{\small\itshape 原收入：银鱼集（后并入《黄裳集·创作卷V》）　｜　年代：待核　｜　地点：—（涉江苏如皋或南京等冒辟疆相关地点，待核）（江南）}\\
写明末四公子之一冒辟疆（号巢民）的斋号'咏怀堂'：追溯冒辟疆与秦淮名妓董小宛的传奇情事（见《影梅庵忆语》）及其不仕清廷的气节；园林、书斋与人格互照，是江南士大夫文化书写的典型个案。注意与E-007《金陵杂记》子目'咏怀堂诗'（1947年南京所见诗集）区分。
\par\vspace{5pt}

\section{黄裳行旅年表}


\begin{quote}\small\itshape
由 \texttt{data/build\_timeline.py} 自动生成，勿手工修改。 仅含 \texttt{approximate\_year} 有值的篇目；年份均待核，「确定」者另见 \texttt{data/huangshang\_essays.json}。 图例：★ 必收　○ 可选　· 资料库
\end{quote}

\subsection{1942}


\begin{itemize}
\item ★ 【辑3】\textbf{白门秋柳} — 南京（白门）
\end{itemize}


\subsection{约1944—1946}


\begin{itemize}
\item ★ 【辑2】\textbf{印度小夜曲} — 印度（加尔各答、菩提加雅等）（待核）
\item ★ 【辑1】\textbf{宝鸡——广元} — 宝鸡（陕）—广元（川），蜀道（待核）
\item ○ 【辑1】\textbf{过徐州} — 徐州（待核）
\item ★ 【辑1】\textbf{成都散记} — 成都（待核）
\item ★ 【辑1】\textbf{昆明杂记} — 昆明（待核）
\item ○ 【辑1】\textbf{森林·雨季·山头人} — 西南/滇缅印一带（待核）（待核）
\item ○ 【辑1】\textbf{音尘} — 西南后方（待核）（推测）
\item ○ 【辑1】\textbf{茶馆} — 西南（成渝一带，待核）（待核）
\item ○ 【辑1】\textbf{江上杂记} — 长江（舟行，具体江段待核）（待核）
\item ○ 【辑1】\textbf{断片} — （追忆，非行旅地点）（待核）
\end{itemize}


\subsection{约1944—1947}


\begin{itemize}
\item ○ 【辑2】\textbf{前线景象} — 前线（待核，疑滇西/缅北）（待核）
\item ★ 【辑2】\textbf{他们怎样看中国} — 西南后方（待核）（待核）
\item · 【辑2】\textbf{“中国通”} — 西南后方（待核）（待核）
\item ○ 【辑2】\textbf{人种·职业·人性的大集合} — 西南后方（待核）（待核）
\item ○ 【辑2】\textbf{几个人物} — 西南后方（待核）（待核）
\item ★ 【辑2】\textbf{中国将军怎样应付美国兵} — 西南后方（待核）（待核）
\item ★ 【辑2】\textbf{美国兵与女人} — 西南后方（待核）（待核）
\item · 【辑2】\textbf{“伟大”的S.O.S.} — 西南后方（待核）（待核）
\item · 【辑2】\textbf{咖啡与战斗力} — 西南后方（待核）（待核）
\item ○ 【辑2】\textbf{种种惊异} — 西南后方（待核）（待核）
\item · 【辑2】\textbf{军中文化} — 西南后方（待核）（待核）
\item ★ 【辑2】\textbf{为美国兵活着的人们} — 西南后方（待核）（待核）
\item ★ 【辑2】\textbf{关于“翻译官”} — 西南后方（待核）（待核）
\item · 【辑2】\textbf{汽车团} — 西南后方（待核，疑滇缅公路沿线）（待核）
\end{itemize}


\subsection{1946}


\begin{itemize}
\item ○ 【辑3】\textbf{旅京随笔} — 南京（待核）
\item ○ 【辑3】\textbf{关于“泽存书库”} — 南京（待核）
\item ○ 【辑3】\textbf{访“盔山精舍”} — 南京（待核）
\item ○ 【辑3】\textbf{「美人肝」} — 南京（待核）
\item ★ 【辑3】\textbf{老虎桥边看“知堂”} — 南京·老虎桥监狱（待核）
\item ★ 【辑3】\textbf{鸡鸣寺} — 南京·鸡鸣寺（待核）
\end{itemize}


\subsection{1947}


\begin{itemize}
\item ★ 【辑3】\textbf{金陵杂记} — 南京（多处）（待核）
\end{itemize}


\subsection{约1947}


\begin{itemize}
\item · 【辑2】\textbf{叙言} — —（待核）
\end{itemize}


\subsection{1949}


\begin{itemize}
\item ○ 【辑3】\textbf{司徒雷登夹着皮包走了以后} — 南京（待核）
\item ○ 【辑3】\textbf{社会科学者眼里的南京} — 南京（待核）
\item ○ 【辑3】\textbf{不再做花瓶——记中国农业机械公司南京工厂} — 南京（待核）
\item ★ 【辑3】\textbf{解放后看江南} — 江南／南京（待核）
\end{itemize}


\subsection{1979}


\begin{itemize}
\item ○ 【辑3】\textbf{南唐二陵} — 南京·南唐二陵（待核）
\item ★ 【辑3】\textbf{天王府的西花园} — 南京·天王府（太平天国旧址）（待核）
\item ○ 【辑3】\textbf{扫叶楼} — 南京·扫叶楼（清凉山）（待核）
\item ○ 【辑3】\textbf{梅园（1979版）} — 南京·梅园（待核）
\item ○ 【辑3】\textbf{白鹭洲公园补记} — 南京·白鹭洲（待核）
\item ○ 【辑3】\textbf{石巢园} — 南京·石巢园（待核）
\item ★ 【辑3】\textbf{秦淮拾梦记} — 南京·秦淮（待核）
\item ○ 【辑3】\textbf{莫愁湖（1979版）} — 南京·莫愁湖（待核）
\item ★ 【辑3】\textbf{重过鸡鸣寺} — 南京·鸡鸣寺（待核）
\item ○ 【辑3】\textbf{献花岩} — 南京郊外·献花岩（具体地点待核）（待核）
\item ★ 【辑3】\textbf{王介甫与金陵} — 南京（半山园等王安石相关遗迹）（待核）
\item ○ 【辑3】\textbf{白下书简} — 南京（白下）（待核）
\end{itemize}


\subsection{年份待核}


\begin{itemize}
\item ★ 【辑8】\textbf{余淡心与金陵} — 南京（秦淮）
\item ○ 【辑8】\textbf{关于余淡心} — —（涉南京秦淮等余怀相关地点）
\item ○ 【辑8】\textbf{关于吴梅村} — —（涉江南，吴伟业相关行迹）
\item ○ 【辑8】\textbf{咏怀堂} — —（涉江苏如皋或南京等冒辟疆相关地点，待核）
\item ○ 【辑4】\textbf{杭州杂记} — 杭州
\item ○ 【辑4】\textbf{海滨消夏记} — 海滨（具体地待核，疑青岛/北戴河一带）
\item ★ 【辑5】\textbf{琉璃厂} — 北京·琉璃厂
\item ○ 【辑8】\textbf{绝代的散文家——张宗子} — —（涉绍兴/杭州等张岱行迹）
\item ★ 【辑7】\textbf{西泠访书记} — 杭州·西泠（西湖孤山一带）
\item ○ 【辑4】\textbf{采石·当涂·青山} — 采石矶、当涂、青山（皖南／长江沿线）
\end{itemize}


\medskip\hrule\medskip

\subsection{各辑篇目统计}


总计 59 篇


\begin{itemize}
\item 第1辑 蜀道与战时中国：9 篇
\item 第2辑 美国兵与异国战争经验：16 篇
\item 第3辑 金陵旧梦：24 篇
\item 第4辑 江南行旅：3 篇
\item 第5辑 北平与旧都余韵：1 篇
\item 第7辑 书肆、旧书与纸上行旅：1 篇
\item 第8辑 人物途中：5 篇
\end{itemize}


\section{资料来源与版本目录（Source Bibliography）}


\begin{quote}\small\itshape
标注约定：\textbf{确定} = 至少一处可靠来源印证；\textbf{待核} = 仅得二手/网络信息，须以实物或图书馆书目核定；\textbf{推测} = 据常识或篇名推断。 本目录记录「我们依据什么」，并区分\textbf{原始结集}与\textbf{后出重印/全集本}。版权未失效，使用以索引与核校为限。
\end{quote}

\medskip\hrule\medskip

\subsection{一、作者基本信息}


\begin{itemize}
\item 黄裳，原名\textbf{容鼎昌}，1919—2012，祖籍山东益都，\textbf{确定}。生于河北井陉（\textbf{待核}）。
\item 散文家、记者、藏书家、古籍版本学者、翻译家（译屠格涅夫《猎人日记/笔记》等）。\textbf{确定}。
\item 1943年应征入伍；1944—1946年在成都、重庆、昆明、桂林、贵阳及印度等地任\textbf{美军译员}——此为《关于美国兵》与《锦帆集》西南诸篇的直接经验背景。\textbf{确定}。
\item 长期任职上海《文汇报》。\textbf{确定}。
\item 2012年9月5日在上海逝世，享年93岁。\textbf{确定}。
\end{itemize}


\begin{quote}\small\itshape
注意：部分网络条目把卒年误作2008；以2012为准（多家讣闻、维基、观察者网报道一致）。
\end{quote}

\medskip\hrule\medskip

\subsection{二、原始结集（民国—1980年代）}


\begin{center}
\tiny
\begin{tabularx}{\linewidth}{@{}XXXXXX@{}}
\toprule
\textbf{书名} & \textbf{出版社} & \textbf{年份} & \textbf{丛书} & \textbf{置信度} & \textbf{备注} \\
\midrule
锦帆集 & 中华书局 & 1946 & 中华文艺丛刊（待核） & 年份\textbf{确定}（多源） & 第一部散文集；书名取李商隐「锦帆应是到天涯」 \\
锦帆集外 & 文化生活出版社 & 1948 & 巴金主编「文学丛刊」第九集（待核） & \textbf{待核（多源）} & 《锦帆集》集外文字 \\
关于美国兵 & 上海出版公司 & 1947 & 周报丛书第一辑 & \textbf{待核（多源）} & 已得完整15篇目录（见第五节B） \\
金陵五记 & 江苏人民出版社 & 1982 & — & \textbf{待核} & 含1942/1946/1947/1949/1979 五时间层；另有《金陵杂记》（金陵书画社，1982）为节取本 \\
花步集 & 生活·读书·新知三联书店 & 1982 & — & \textbf{待核（单源）} & 后并入《黄裳集·创作卷IV》 \\
晚春的行旅 & 待核（疑山西人民，未证实） & 待核 & — & \textbf{待核} & 书名直指行旅主题；后并入《黄裳集·创作卷IV》 \\
\bottomrule
\end{tabularx}
\end{center}


\begin{quote}\small\itshape
\textbf{ 重要更正——《白门秋柳》不是民国原集}：豆瓣 subject 26801544 的《白门秋柳》是\textbf{后出主题选本}，分「秦淮拾梦／闲情逸事／驻足闲览／羁旅山川／榆下梦痕／月旦人物」六辑，汇黄裳数十年散文精华。它与1942年\textbf{单篇}《白门秋柳》同名。该单篇初出《锦帆集》（1946），后又收入《金陵五记》。本项目据此把「白门秋柳」从「原始结集」改列：单篇归《锦帆集》，同名选本移入「现有市面选本」（见第四节）作对照对象。
\end{quote}

\subsubsection{其他相关集（书话/记人/记戏/古籍，部分已得目录）}

\begin{itemize}
\item 《榆下说书》《银鱼集》（《黄裳集·创作卷V》）——书话/记人，第七、八辑主要来源。\textbf{已得目录线索}（见第五节D）。
\item 《翠墨集》——安徽教育出版社，2006（"黄裳作品系列"，264页，\textbf{待核}）；后并入《黄裳集·创作卷VI》。内容含"清雅书话、珍本题跋、\textbf{山川风物}、历史人事"，是第四/六/八辑可采之源；\textbf{逐篇目录未得}。
\item 《榆下杂说》（《黄裳集·创作卷VI》）——书话/杂说，第六、七辑补采。
\item 《旧戏新谈》——1947—1948年在《文汇报》副刊所写戏剧杂谈，\textbf{开明书店1948年结集}（约50余篇）；后有"大家小书"系列（北京出版社，2003）。\textbf{性质提示：整体为剧评/戏考，多无行旅线索，原则上不直接入本行旅选本}；仅其中"剧场现场""伶人记"类带地点/人物者可酌选入第八辑（戏曲人物）或第五辑（剧场）。已见篇名如《春闺梦》《一捧雪》等均属戏目评谈，暂不收。
\item 《珠还记幸》《来燕榭文存》（《黄裳集·创作卷XIV》）《前尘梦影新录》——记人、版本、旧事来源，供第五/六/七/八辑补采。\textbf{均待核}。
\item 《笔祸史谈丛》《妆台杂记》（《黄裳集·创作卷IX》）——史谈/杂记，备查。
\end{itemize}


\medskip\hrule\medskip

\subsection{三、现代全集 / 底本（核对基准）}


\begin{center}
\scriptsize
\begin{tabularx}{\linewidth}{@{}XXXX@{}}
\toprule
\textbf{名称} & \textbf{出版社} & \textbf{说明} & \textbf{置信度} \\
\midrule
\textbf{黄裳集} & 山东人民出版社 & 目前最全；分\textbf{创作卷（约16卷）+ 古籍研究卷 + 译文卷 + 书信卷}，大致按出版时间分类编入。\textbf{建议作为主要底本。}《锦帆集》《锦帆集外》《关于美国兵》在创作卷I；《花步集》《金陵五记》《晚春的行旅》在创作卷IV。 & 结构\textbf{确定}；卷次/年份\textbf{待核} \\
来燕榭少作五种 & 中华书局 & 早期作品五种合刊，可校早期文本 & \textbf{待核} \\
黄裳文集（6册） & — & 早出文集，可作旁证 & \textbf{待核} \\
\bottomrule
\end{tabularx}
\end{center}


\medskip\hrule\medskip

\subsection{四、现有市面选本（拟对照与超越的对象）}


\begin{quote}\small\itshape
用于第三阶段「现有选本不足」分析。各本编选思路、收文范围、注释水平\textbf{待逐一翻阅核对}。
\end{quote}

\begin{itemize}
\item 《黄裳散文》（多家版本）——通代散文选，非行旅专题。
\item 《黄裳自选集》——作者自选，分辑，含早期西南/美国兵诸篇（如昆明杂记、茶馆、森林·雨季·山头人、美国兵与女人等），是核早期篇目的有用线索。
\item 《山川、历史、人物》——主题向选本，与本项目取向最接近，须重点对照。
\item 商务印书馆「流金文丛」重印《金陵五记》等——便于核校金陵部分文本。
\end{itemize}


\medskip\hrule\medskip

\subsection{五、各书已得目录（来源：网络检索的目录，\textbf{均待与实物核对}）}


\subsubsection{A.《金陵五记》编年目录}


\begin{itemize}
\item \textbf{1942}：白门秋柳
\item \textbf{1946}：旅京随笔；鸡鸣寺；关于「泽存书库」；访「盔山精舍」；「美人肝」；老虎桥边看「知堂」
\item \textbf{1947}：金陵杂记（小序＋半山寺与谢公墩、石观音寺、周处读书台、绛云书卷美人图、柳如是、快园、随园、梅园、后湖、鸡鹅巷与裤子裆、小虹桥、咏怀堂诗、梅花山、莫愁湖、明太祖与徐达、桃花扇底看南朝、马湘兰、记者生涯、豁蒙楼上看浓春、燕子矶、白鹭洲）
\item \textbf{1949}：解放后看江南；社会科学者眼里的南京；司徒雷登夹着皮包走了以后；不再做花瓶——记中国农业机械公司南京工厂
\item \textbf{1979}：白下书简；秦淮拾梦记；白鹭洲公园补记；重过鸡鸣寺；天王府的西花园；梅园；王介甫与金陵；莫愁湖；石巢园；扫叶楼；献花岩；南唐二陵
\item \textbf{附录}：余怀诗两种（《咏怀古迹》《味外轩集》）
\end{itemize}


\begin{quote}\small\itshape
上列子目个别字（如「莫愁湖/奠愁湖」）疑为录入异写，须以原书校正。
\end{quote}

\subsubsection{B.《关于美国兵》目录（15篇，按原书次序）}

叙言；他们怎样看中国；"中国通"；人种·职业·人性的大集合；几个人物；中国将军怎样应付美国兵；美国兵与女人；"伟大"的S.O.S.；前线景象；咖啡与战斗力；种种惊异；军中文化；汽车团；为美国兵活着的人们；关于"翻译官"。

\begin{quote}\small\itshape
此为第二辑完整骨架。各篇确切写作时间/地点（多笼统为西南后方或前线）待据原文定位；"S.O.S."疑指美军后勤勤务部（Services of Supply），待考。
\end{quote}

\subsubsection{C.《锦帆集》目录（部分）}

断片；白门秋柳；过徐州；宝鸡——广元；成都散记；音尘；江上杂记；"江湖"后记；后记。

\begin{quote}\small\itshape
含蜀道（宝鸡——广元）、蜀中（成都散记）、江北（过徐州）、江河（江上杂记）诸篇，为第一辑骨架。
\end{quote}

\subsubsection{D.《榆下说书》《银鱼集》目录（部分）}

\begin{itemize}
\item \textbf{榆下说书}：书的故事；古书的作伪；西泠访书记；谈"善本"；谈"题跋"；谈"集部"；谈禁书；书痴 等。
\item \textbf{银鱼集}：绝代的散文家——张宗子；关于张宗子；关于余淡心；余淡心与金陵；关于吴梅村；咏怀堂；海滨消夏记 等。
\end{itemize}

\begin{quote}\small\itshape
取舍提示：\textbf{西泠访书记}（杭州访书，带行踪）入第七辑；\textbf{余淡心与金陵 / 张宗子 / 吴梅村}等记人之作入第八辑；\textbf{谈"善本"/题跋/集部/禁书 / 古书的作伪}等纯书话无行踪者\textbf{不入本选本}。
\end{quote}

\medskip\hrule\medskip

\subsection{六、检索线索（供继续查证）}


\begin{itemize}
\item 图书馆书目：中国国家图书馆、上海图书馆、南京图书馆 OPAC（核原始版权页最权威）。
\item 百科/条目：维基百科「黄裳」「容鼎昌」、《中国大百科全书》第三版、百度百科（百度百科 WebFetch 曾返回 403，须改用浏览器或镜像）。
\item 书目站：豆瓣读书「黄裳」著者页与「黄裳集」丛书页、Google Books。
\item 报刊副刊：《人民日报》海外版、《中华读书报》《新民晚报》相关纪念/评介文章。
\end{itemize}


\begin{quote}\small\itshape
所有上述均为\textbf{二手}，凡涉及精确年份/出版社/篇名，须以\textbf{原书或图书馆书目}为最终依据。
\end{quote}

\section{收录书目（书锚）与版本}
\noindent\small 本项目据以采辑的原始结集与全集本；版权均未失效，仅作版本核校之用。\normalsize\medskip
\noindent\textbf{《锦帆集》}\,{\small（中华书局，1946，中华文艺丛刊）}\\
作者第一部散文集；书名取李商隐（义山）'锦帆应是到天涯'句。丛刊归属已多源印证（确定）。章序：ch1序、ch2断片、ch3白门秋柳、ch4过徐州、ch5宝鸡——广元、ch6成都散记、ch7音尘、ch8江上杂记、ch9'江湖'后记、ch10后记。多记抗战后期西南/蜀道行旅。注意：昆明杂记不在此集，在《锦帆集外》。
\\{\small\itshape 来源：多源印证（豆瓣'锦帆集 锦帆集外 关于美国兵'条目、人民日报副刊、百度百科）；丛刊名'中华文艺丛刊'多源确认。}
\par\vspace{5pt}
\noindent\textbf{《锦帆集外》}\,{\small（文化生活出版社，1948，文学丛刊第九集（巴金主编，多源印证））}\\
《锦帆集》集外文字结集，列巴金主持的'文学丛刊第九集'（多源印证）。核心篇目：《昆明杂记》（五部分；E-018）。完整篇目尚待录入。
\\{\small\itshape 来源：网络检索（豆瓣/微信读书；多处提及巴金文学丛刊第九集）；昆明杂记在此集已多源印证。}
\par\vspace{5pt}
\noindent\textbf{《音尘集》}\,{\small（辽宁教育出版社，1996，书趣文丛第三辑）}\\
含三部分合订：（1）《锦帆集》重印；（2）《印度小夜曲》（新增，8章：序曲/雨季/加尔各答/菩提加雅/群莺乱飞/弄蛇者/女人们/森林雨季；写1944—1946年印度行役）；（3）《关于美国兵》重印。《印度小夜曲》是目前能确认的印度行旅专题长篇，为本项目重要新增来源。
\\{\small\itshape 来源：网络检索（多源提及'音尘集'辽宁教育版，书趣文丛第三辑，1996）；年份与出版社待实物核。}
\par\vspace{5pt}
\noindent\textbf{《关于美国兵》}\,{\small（上海出版公司，1947，周报丛书第一辑）}\\
作者抗战末随美军任译员期间的现场观察。已得完整目录（15篇，见 essays 中 E-020..E-034）。
\\{\small\itshape 来源：多源印证（豆瓣条目、综述文章）。}
\par\vspace{5pt}
\noindent\textbf{《金陵五记》}\,{\small（金陵书画社，1982）}\\
汇集1942—1979年写南京文字，分五个时间层。1982年6月初版，260页，统一书号10234-004，定价¥0.85（金陵书画社版，已多源确认）。另有后出江苏古籍出版社版（ISBN 9787806431757）。旧版录'江苏人民出版社'为误，已更正。全书核心来源。
\\{\small\itshape 来源：网络检索（百度百科'金陵五记'、微信读书/中图网目录、商务'流金文丛'重印页）；金陵书画社1982版书号统一书号10234-004，多源确认。}
\par\vspace{5pt}
\noindent\textbf{《花步集》}\,{\small（生活·读书·新知三联书店，1982）}\\
后并入《黄裳集·创作卷IV》（与《金陵五记》《晚春的行旅》同册）。篇目待录入。
\\{\small\itshape 来源：网络检索（图书馆书目、豆瓣《黄裳集》系列）；出版社/年份单源，待核。}
\par\vspace{5pt}
\noindent\textbf{《晚春的行旅》}\,{\small（待核（疑山西人民出版社，未证实），待核（疑1980年代初））}\\
书名直指'行旅'主题，应为本选本重点采辑对象。后并入《黄裳集·创作卷IV》。篇目与出版信息均待核。
\\{\small\itshape 来源：网络检索（豆瓣《黄裳集》系列）；出版社/年份未坐实。}
\par\vspace{5pt}
\noindent\textbf{《榆下说书 / 银鱼集》}\,{\small（（原版年份待核）后并入《黄裳集·创作卷V》（山东人民出版社），待核）}\\
书话集。《榆下说书》得目录线索：书的故事、古书的作伪、西泠访书记、谈'善本'、谈'题跋'、谈'集部'、谈禁书、书痴等；《银鱼集》：绝代的散文家——张宗子、关于张宗子、关于余淡心、余淡心与金陵、关于吴梅村、咏怀堂、海滨消夏记等。为第七、八辑主要来源。
\\{\small\itshape 来源：网络检索（中图网/苏宁《黄裳集·创作卷V》页、Google Books）。}
\par\vspace{5pt}
\noindent\textbf{《翠墨集 / 榆下杂说》}\,{\small（后并入《黄裳集·创作卷VI》（山东人民出版社），待核）}\\
书话/杂说集，供第六、七辑补采。篇目待录入。
\\{\small\itshape 来源：网络检索（苏宁《黄裳集·创作卷VI》页）。}
\par\vspace{5pt}
\noindent\textbf{《黄裳集（山东人民出版社）》}\,{\small（山东人民出版社，约2017—2022（分批））}\\
目前最全的作品集，分创作卷（约16卷）+古籍研究卷+译文卷+书信卷。建议作为主要底本与核对基准。创作卷I=锦帆集/锦帆集外/关于美国兵；IV=花步集/金陵五记/晚春的行旅；V=榆下说书/银鱼集；VI=翠墨集/榆下杂说。仍在版权保护期内。
\\{\small\itshape 来源：网络检索（豆瓣'黄裳集'丛书页、中图网团购页、苏宁分卷页）。卷次/年份待核。}
\par\vspace{5pt}

\section{待核清单与已知缺口}
\subsection{待核清单（needs\_verification）}
\begin{itemize}
\item 【已更新】《锦帆集》丛刊归属已多源确认为'中华文艺丛刊'；《锦帆集外》已多源确认为'文学丛刊第九集（巴金主编）'；仍待图书馆书目最终确认。
\item 【已更新】《金陵五记》出版社已更正为'金陵书画社'（1982年6月，统一书号10234-004）；旧录'江苏人民出版社'系误，已删；仍待实物核定。《花步集》《晚春的行旅》出版社/年份仍未坐实。
\item 【已更新】《昆明杂记》已更正归属《锦帆集外》（非《锦帆集》）；五部分结构已多源确认；仍待原书核定。
\item 【已更新】《音尘集》（辽宁教育出版社，1996，书趣文丛第三辑）及《印度小夜曲》已录入；年份、出版社待实物核。
\item 《白门秋柳》单篇（1942，初出《锦帆集》ch3）与同名后出选本（豆瓣26801544）须严格区分。白门秋柳是否存在'锦帆集本'与'金陵五记本'异文，待原书比对。
\item 关于美国兵各篇的确切写作时间与发生地点（多笼统标'西南后方/前线'，待据原文定位）。
\item 《重过鸡鸣寺》（E-012）写作年份：《金陵五记》编年列1979，一说1980；待以原书内文日期核定。
\item 西泠访书记、琉璃厂、海滨消夏记的写作年份与原书归属仍待核。
\item 《关于美国兵》15篇个别字词（如'S.O.S.'所指建制）须以原书校。
\item 印度小夜曲八章内容须以《音尘集》原书核定（八章章名已多源确认，内容细节待核）。
\end{itemize}
\subsection{已知缺口（gaps）}
\begin{itemize}
\item 第三辑（金陵）：《金陵五记》五个时间层已基本录入（59篇中24篇属第三辑）；梅园(1979)与莫愁湖(1979)须核定是否与金陵杂记子目重合；「美人肝」内容完全待核。
\item 《锦帆集外》：《昆明杂记》（五部分）已录入（E-018）；其余集外文字（除《昆明杂记》外）完整篇目尚待录入。
\item 第二辑新增：《印度小夜曲》（E-060，《音尘集》1996）已录入为'必收'；八章结构确认，内容细节待原书核定。
\item 第四辑（江南）：仍仅3篇，《花步集》《晚春的行旅》逐篇目录获取受阻（豆瓣/中图网 403），须改走图书馆 OPAC（国图/上图）或《黄裳集·创作卷IV》原书。
\item 第五辑（北平/琉璃厂/旧城/剧场）仅《琉璃厂》1篇，须从《翠墨集》《珠还记幸》《来燕榭文存》补采旧城/剧场/访书诸篇。
\item 第六辑（古迹/废墟，跨地域）独立篇目为零；天王府等怀古篇因金陵地缘划入第三辑；须从《翠墨集》'山川风物'内容（目录待录入）补充跨地域古迹篇。
\item 戏曲题材（剧场/戏曲人物）篇目几乎空白，须查《旧戏新谈》中含地点/人物的篇目，选入第五辑或第八辑。
\item 第七辑（书旅）仅《西泠访书记》1篇；须从《翠墨集》《榆下杂说》《珠还记幸》中摘选带行踪的访书文。
\item 第八辑（人物）现5篇；《珠还记幸》《来燕榭文存》中记人文字尚未录入。
\item 多部原始结集的出版社/年份仍为网络二手，须以图书馆书目/实物核定（金陵五记出版社已更正，其他仍待核）。
\item 尚未系统排查'同篇异文'（如白门秋柳的锦帆集本 vs 金陵五记本；梅园/莫愁湖1947 vs 1979版）。
\end{itemize}
\subsection{不宜收判例（not\_recommended）}
\begin{itemize}
\item \textbf{纯书话（无行踪）：谈"善本"、谈"题跋"、谈"集部"、谈禁书、古书的作伪 等（《榆下说书》）}　属版本学/书话专论，无地点或行旅线索，与'行旅文史'主旨偏离；留待黄裳书话专集，不入本选本。
\item \textbf{纯文学评论：关于张宗子 / 关于吴梅村 / 关于余淡心 等单纯论人论文之作}　论人论文成分重、行旅成分弱；除非与具体地点（如余淡心与金陵）绑定，否则列'不宜收'或'可选'，避免冲淡行旅主线。
\item \textbf{纯时政评论篇（如部分1949年时评）}　时评性强、时效性内容多，须经史地注释方能入；超出'地点触发的文史随笔'范畴者不收。
\end{itemize}

\end{document}

